\subsection{预备定理}

在本节中，A 表示一个局部环，m 表示其极大理想，k = A/m 表示其剩余域。设 A 关于 m-进拓扑是 Hausdorff 且完备的。令 B = A{[}{[}X{]}{]}；它是一个局部环，其极大理想 N 由 m 和 X 生成；关于 \(\mathfrak{N}\)-进拓扑，B 是 Hausdorff 且完备的（第三章，§ 2，第 6 节，命题 6）。

对于每个形式幂级数

\[
f = \sum_ {i = 0} ^ {\infty} a _ {i} \mathbf {X} ^ {i} \in \mathbf {B},
\]

记作

\[
\bar {f} = \sum_ {i = 0} ^ {\infty} \bar {a} _ {i} X ^ {i} \in k [ [ X ] ],
\]

其中 \(\bar{a}_{i}\) 表示 \(a_{i}\) 在 k 中的典范像。级数 \(\bar{f}\) 称为 f 的约化幂级数；若 \(\bar{f} \neq 0\)，则 \(\bar{f}\) 的阶（即满足 \(a_{s} \notin m\) 的最小整数 s）称为 f 的约化阶。

命题 5. 设 \(f \in \mathbf{B}\) 是约化幂级数非零的级数。令 \(s\) 表示其约化阶，令 \(\mathbf{M}\) 为 \(\mathbf{B}\) 的子-A-模，其基为 \(\{1, X, \ldots, X^{s-1}\}\)。则 \(\mathbf{B}\) 是 \(\mathbf{M}\) 与 \(f\mathbf{B}\) 的直和，且 \(f\) 不是 \(\mathbf{B}\) 中的零因子。

\begin{enumerate}
\def\labelenumi{(\alph{enumi})}
\tightlist
\item
  我们证明 \(f\mathbf{B} \cap \mathbf{M} = (0)\)。假设存在一个关系：
\end{enumerate}

\[
\left(\sum_ {i = 0} ^ {\infty} b _ {i} \mathrm{X} ^ {i}\right) f = r _ {0} + r _ {1} \mathrm{X} + \dots + r _ {s - 1} \mathrm{X} ^ {s - 1} \quad (b _ {i} \in \mathrm{A}, r _ {j} \in \mathrm{A}).\tag{3}
\]

我们证明所有 \(b_{i}\)（从而所有 \(r_{j}\)）都为零；这尤其将证明 \(f\) 不是 B 中的零因子。由于 A 是 Hausdorff，只需证明对所有 \(b_{i} \in \mathfrak{m}^{n}\)\(i \geqslant 0\) 和所有 \(n \geqslant 0\)，都有 。\(n = 0\) 时显然成立。我们作双重归纳：假设对所有  有 \(b_{i} \in \mathfrak{m}^{n-1}\)，且对 \(i\) 有 \(b_{i} \in \mathfrak{m}^{n}\)，并证明这推出 \(i < k\)\(b_{k} \in \mathfrak{m}^{n}\)。为此，写成 \(f = \sum_{i=1}^{\infty} a_{i} X^{i}\)，并在（3）中比较 \(\mathbf{X}^{s+k}\) 的系数；于是：

\[
\left(b _ {0} a _ {s + k} + \dots + b _ {k - 1} a _ {s + 1}\right) + b _ {k} a _ {s} + \left(b _ {k + 1} a _ {s - 1} + \dots + b _ {k + s} a _ {0}\right) = 0.\tag{4}
\]

第一括号中的项属于 \(\mathfrak{m}^n\)，因为当  时 \(b_i \in \mathfrak{m}^n\)；第二括号中的项同理，因为对所有 \(i < k\) 都有 \(b_i \in \mathfrak{m}^{n-1}\)，且当 \(i\) 时 \(a_i \in \mathfrak{m}\)。因此 \(i \leqslant s - 1\)\(b_k a_s \in \mathfrak{m}^n\)，又由于 \(a_s\) 是 A 中的可逆元，故 \(b_k \in \mathfrak{m}^n\)。

\begin{enumerate}
\def\labelenumi{(\alph{enumi})}
\setcounter{enumi}{1}
\tightlist
\item
  我们证明 \(f\mathbf{B} + \mathbf{M} = \mathbf{B}\)。写成
\end{enumerate}

\[
g = a _ {s} + a _ {s + 1} \mathbf {X} + a _ {s + 2} \mathbf {X} ^ {2} + \dots ;
\]

它是 B 中的可逆元。于是

\[
f - \mathrm{X} ^ {s} g = a _ {0} + a _ {1} \mathrm{X} + \dots + a _ {s - 1} \mathrm{X} ^ {s - 1};
\]

因此若写成 \(fg^{-1} - \mathbf{X}^s = (f - \mathbf{X}^s g)g^{-1} = -h\)，则 h 的系数属于 \(h\)\(m\)。令 r 为 B 的一个元素。按 n 作归纳，定义 B 中元素的序列 \(r\)\(n\)\((q^{(n)})\)：取 \(q^{(0)}\) 为满足下式的唯一级数：

\[
r \equiv \mathbf {X} ^ {s} q ^ {(0)} \pmod {\mathbf {M}};\tag{5}
\]

写成 \(h = \sum_{i=0}^{\infty} h_i X^i\) 和 \(q^{(n)} = \sum_{i=0}^{\infty} q_i^{(n)} X^i\)，则 \(q_i^{(n)}\) 由下式定义：

\[
q _ {i} ^ {(n)} = \sum_ {j = 0} ^ {i + s} h _ {j} q _ {i + s - j} ^ {(n - 1)}.\tag{6}
\]

由（6）立即得到：

\[
\mathbf {X} ^ {s} q ^ {(n)} \equiv h q ^ {(n - 1)} \pmod {\mathbf {M}}.\tag{7}
\]

由于对所有  都有 \(h_j \in \mathfrak{m}\)，由（6）对 n 作归纳还可得到，对所有 \(j\)\(n\) 都有 \(q_i^{(n)} \in \mathfrak{m}^n\)。由于 A 是完备的，级数\(i\)\(n\)\(A\)

\[
q ^ {(0)} + q ^ {(1)} + \dots + q ^ {(n)} + \dots
\]

收敛到 B 中的一个元素 q。由（5）和（7），\(q\)

\[
\mathbf {X} ^ {s} (q ^ {(0)} + q ^ {(1)} + \dots + q ^ {(n)}) \equiv r + h (q ^ {(0)} + \dots + q ^ {(n - 1)}) \pmod {\mathbf {M}}. \tag {8}
\]

由于 \(\mathbf{M}\) 是闭的，（8）在极限处给出 \(r \equiv (\mathbf{X}^s - h)q\)（模 \(\mathbf{M}\)），即

\[
r \in f g ^ {- 1} q + \mathbf {M} \subset f \mathbf {B} + \mathbf {M}.
\]

也可以利用第三章 §2 的结果证明关系 \(B = fB + M\)（参见习题 12）。这里采用的方法的优点是也适用于收敛级数。

推论。在命题 5 的假设和记号下，设 \(s \geqslant 1\)，因而 \(f \in \mathbf{B}\mathfrak{m} + \mathbf{B}\mathbf{X}\)。则从 \(h\)\(\mathbf{B}' = \mathbf{A}[[\mathbf{T}]]\) 到 \(\mathbf{B} = \mathbf{A}[[\mathbf{X}]]\) 的 A-同态 ，满足 \(h(\mathbf{T}) = f\)（第三章，§2，第 9 节，命题 11 (a)），在 \(\mathbf{B}\) 上定义一个以  为基的自由 \(\mathbf{B}'\)-模结构。特别地，\(\{1, \mathbf{X}, \ldots, \mathbf{X}^{s-1}\}\)\(h\) 是单射。

赋予 B' 模 B 以 (T)-进滤过，该滤过由  时的 \(f^{n}B\) 组成（第三章，§2，第 1 节）。则 B/fB 是环 A = B'/TB' 上的自由模，而 \(n \geqslant 0\)\(X^{i}\)（\(0 \leqslant i \leqslant s - 1\)）在此 A-模中的像构成其基（命题 5）；又由于 f 不是 B 中的零因子（命题 5），\(Bf^{n}/Bf^{n+1}\) 也是秩为 s 的 (B'/TB')-自由模，因而第三章 §2 第 8 节的条件 (GR) 得到满足（以 B' 代替 A，以 B 代替 M）。另一方面，由于 B' 关于 (T)-进滤过是 Hausdorff 且完备的，且由上可知 gr(B) 是有限生成的 gr(B')-模，首先可见（第三章，§2，第 9 节，命题 12 的推论 1）B 是有限生成的 B' 模。于是推论的第一断言由第三章 §2 第 9 节命题 13 得出。第二断言立即由此得出。

定义 2. 若多项式 \(\mathbf{F} \in \mathbf{A}[\mathbf{X}]\) 形如

\[
\mathbf {F} = \mathbf {X} ^ {s} + a _ {s - 1} \mathbf {X} ^ {s - 1} + \dots + a _ {0},
\]

其中 \(a_{i} \in \mathfrak{m}\)（\(0 \leqslant i \leqslant s - 1\)），则称其为特异多项式。

注意，两个特异多项式的乘积仍是特异多项式。

命题 6（预备定理）。设 \(f \in \mathbf{B}\) 是约化幂级数非零的级数，且 \(s\) 是其约化阶。则存在唯一的有序对 \((u, \mathbf{F})\)，其中 \(u\) 是 \(\mathbf{B}\) 中的可逆元，\(\mathbf{F}\) 是次数为 \(s\) 的特异多项式，并且 \(f = u\mathbf{F}\)。

写成 \(F = X^{s} + G\)，其中 \(G = g_{0} + \cdots + g_{s-1}X^{s-1} (g_{i} \in A)\)。关系 f = uF 等价于 \(F = u^{-1}f\)，即 \(X^{s} = u^{-1}f - G\)。因此命题 5 表明 G 和 \(u^{-1}\) 是唯一的，从而 F 和 u 也是唯一的。它还表明存在 \(v \in B\) 和多项式 \(G = g_{0} + \cdots + g_{s-1}X^{s-1} (g_{i} \in A)\)，使得 \(X^{s} = vf - G\)；还需证明 v 在 B 中可逆，且对所有 i 都有 \(g_{i} \in m\)。现在，记 \(\bar{g}_{i}\) 为 \(g_{i}\) 在 k 中的典范像，并以 \(\bar{f}\)、\(\bar{v}\) 表示 f、g 的约化幂级数，

\[
\mathbf {X} ^ {s} + \bar {g} _ {0} + \bar {g} _ {1} \mathbf {X} + \dots + \bar {g} _ {s - 1} \mathbf {X} ^ {s - 1} = \bar {f} \bar {v};
\]

由于 \(\bar{f}\) 的阶为 \(s\)，所以对所有  都有 \(\bar{g}_i = 0\)，且 \(i\)\(\bar{v}\) 的阶为 0，故 v 可逆。\(v\)

命题 7. 设 F 是特异多项式，g、h 是 B 中两个形式幂级数，满足 F = gh。则存在 B 中的可逆元 u，使得 ug 和 \(u^{-1}h\) 是特异多项式，且 \(\mathbf{F} = (ug)(u^{-1}h)\)。

事实上，g 和 h 的约化幂级数均非零；因此由命题 6，存在 B 中的可逆元 u、v，使得 ug 和 vh 是特异多项式。于是 \(\neq0\)\(uvF = (ug)(vh)\) 是特异多项式，且 uv 可逆。考察约化幂级数，立即可见 F 与 uvF 的约化阶相同，即次数相同。因此由命题 6 的唯一性断言，F = uvF，从而 uv = 1。

推论。进一步设 A 是整环，F 是次数为 s 的特异多项式。F 在 A{[}X{]} 中为极端元的充要条件，是它在 B = A{[}{[}X{]}{]} 中为极端元。

假设 F 在 A{[}X{]} 中不是极端元，于是 \(F = f_{1}f_{2}\)，其中 \(f_{1}\) 和 \(f_{2}\) 是 A{[}X{]} 中的不可逆元；由于 \(f_{1}\) 和 \(f_{2}\) 的最高次项系数之积等于 1，这些系数在 A 中可逆，而假设推出 \(f_{1}\) 和 \(f_{2}\) 的次数分别大于 0 且小于 s；由于约化多项式 \textgreater\textless\(\bar{f}_{1}, \bar{f}_{2}\) 满足 \(\bar{f}_{1}\bar{f}_{2} = X^{s}\)，\(\bar{f}_{1}\) 和 \(\bar{f}_{2}\) 都不可能在 k{[}{[}X{]}{]} 中可逆，因为若 \(\bar{f}_{1}\) 可逆，则 \(\bar{f}_{2}\) 的阶将为 s，这是荒谬的。更不用说，\(f_{1}\) 和 \(f_{2}\) 都不在 B 中可逆，F 也不是 B 中的极端元。

反之，若 F 在 A{[}{[}X{]}{]} 中不是极端元，则 F = gh，其中 g、h 在 B 中都不可逆；因此它们的约化阶都 \(\geqslant 1\)；于是命题 7 中的特异多项式 ug 和 \(u^{-1}h\) 都不是常数，这表明 F 不是 A{[}X{]} 中的极端元。

